\documentclass[11pt]{article}
\usepackage[a4paper,margin=26mm]{geometry}
\usepackage{amsmath,amssymb,amsthm,hyperref}
\newtheorem{theorem}{Theorem}
\newtheorem{lemma}[theorem]{Lemma}
\newtheorem{corollary}[theorem]{Corollary}
\newcommand{\F}{\mathbb F}
\newcommand{\GL}{\operatorname{GL}}
\newcommand{\tr}{\operatorname{tr}}
\newcommand{\ord}{\operatorname{ord}}
\title{Two coefficients certify certain prime orders in the natural $G_2$ module}
\author{Abstract preparation by \texttt{/root/mathematical\_review}}
\date{3 November 2026}
\begin{document}
\maketitle

\begin{abstract}
In the natural seven-dimensional module of a split group of type $G_2$,
the characteristic polynomial is determined by its first two elementary
coefficients. If a prime $\ell$ satisfies $\ord_\ell(q)=6$, a finite table
of $(\ell-1)/6$ pairs of coefficients therefore gives an exact test for
order $\ell$ in $G_2(q)$. In odd characteristic one may instead use
$\tr(g)$ and $\tr(g^2)$; the latter needs only twenty-eight products of
matrix entries. The coefficient formulation uses no division and applies
in every characteristic. For $q=11$ we also explain why a word $w$ whose
sixth power has order $37$ must have order $37$ or $111$. This last
argument reuses the previously reviewed involution theorem, and needs no
new assertion about a full torus centralizer. These are abstract proof
and cost statements. No table constants, Original matrix, word or case
have been evaluated, and no implementation or application is adopted.
\end{abstract}

\section{The natural representation and its two coefficients}

Let $k$ be a field, let $\mathbf G$ be the split algebraic group of
type $G_2$, and let $V$ be its natural seven-dimensional module. The
word ``natural'' is essential: it means the usual module obtained from
the trace-zero part of a split Cayley algebra. An identification of an
abstract group with $G_2(q)$, or preservation of a bilinear form alone,
does not identify an arbitrary matrix representation with this module.

We use the standard weight description of this representation. Over an
algebraic closure, a maximal torus has eigenvalues
\begin{equation}\label{eq:weights}
 1,x,y,z,x^{-1},y^{-1},z^{-1},\qquad xyz=1.
\end{equation}
One way to obtain this description is to use the usual $\mathrm{SL}_3$
subgroup in the split Cayley construction. The seven-space restricts as
$1\oplus W\oplus W^*$, where $W$ is its three-dimensional defining
module. Its diagonal torus has $\dim T=2$, the rank of $G_2$, and acts
with the displayed eigenvalues. This is a standard description of the
natural representation, not a recognition theorem for a saved matrix
model. It remains valid in characteristic two; no assertion about a
nondegenerate restricted bilinear form is being made there.

Write $e_j(g)$ for the elementary symmetric functions of the seven
eigenvalues, so that
\[
 \chi_g(X)=X^7-e_1X^6+e_2X^5-e_3X^4+e_4X^3-e_5X^2+e_6X-e_7.
\]
These conventions fix the signs also in small characteristic, where
some of the signs coincide.

\begin{lemma}\label{lem:coefficients}
For every $g\in\mathbf G(k)$, put $t=e_1(g)$ and $E=e_2(g)$. Then
\begin{equation}\label{eq:polynomial}
 \chi_g(X)=X^7-tX^6+EX^5-FX^4+FX^3-EX^2+tX-1,
 \qquad F=t^2+t-E.
\end{equation}
In particular $(e_1,e_2)$ determines the whole characteristic polynomial.
\end{lemma}
\begin{proof}
First suppose that $g$ belongs to the torus in~\eqref{eq:weights}.
Set $s=x+y+z$, $r=xy+yz+zx$, and $A=s+r$. Since $xyz=1$, the six
nonconstant eigenvalues have characteristic polynomial
\begin{align*}
 &(X^3-sX^2+rX-1)(X^3-rX^2+sX-1)\\
 &\quad=X^6-AX^5+(sr+A)X^4-(2+s^2+r^2)X^3
              +(sr+A)X^2-AX+1.
\end{align*}
Multiplication by $X-1$ gives
\[
 e_1=A+1,\qquad e_2=sr+2A,\qquad
 e_3=2+s^2+r^2+sr+A=A^2+A+2-sr.
\]
It follows by polynomial arithmetic over the integers that
\[
 e_3=e_1^2+e_1-e_2.
\]
The inverse pairs in the eigenvalues also give $e_7=1$, $e_6=e_1$,
$e_5=e_2$ and $e_4=e_3$. No division was used, so all these identities
hold in every characteristic.

For completeness, passing from the torus to an arbitrary element uses
the usual Jordan decomposition in a connected reductive algebraic group
and conjugacy of semisimple elements into maximal tori. Extend $k$ to
an algebraic closure and write $g=g_sg_u=g_ug_s$. On each eigenspace of
$g_s$, the commuting factor $g_u$ is unipotent. Thus $g$ and $g_s$ have
the same eigenvalues with multiplicity, and hence the same characteristic
polynomial. Conjugate $g_s$ into the torus and apply the calculation.
The resulting coefficient identity restricts to $\mathbf G(k)$.
\end{proof}

This proof explicitly imports the standard natural weight description,
Jordan decomposition and the maximal torus theorem. It neither proves
an existing Cayley-model membership assertion nor uses a classification
of maximal finite subgroups.

\section{A finite table for a prime of extension degree six}

Let $q$ be a prime power and let $\ell$ be a prime with
$\ord_\ell(q)=6$. In particular $\ell$ is different from the
characteristic. Choose a primitive $\ell$th root $\zeta$ in an
algebraic closure of $\F_q$, and put
\[
 H=\langle q\rangle\leq(\mathbb Z/\ell\mathbb Z)^\times.
\]
It has order six. For each coset $kH$ define
\begin{equation}\label{eq:table}
 f_k(X)=\prod_{i=0}^5(X-\zeta^{kq^i}),\qquad
 P_k(X)=(X-1)f_k(X),\qquad
 \sigma_k=(e_1(P_k),e_2(P_k))\in\F_q^2.
\end{equation}
Here the elementary coefficients refer to the seven roots of $P_k$,
with the sign convention of the preceding section. The polynomials
$f_k$ are exactly the distinct irreducible factors of $\Phi_\ell$
over $\F_q$.

\begin{lemma}\label{lem:possible}
An element of $\GL_7(\F_q)$ has exact order $\ell$ if and only if its
characteristic polynomial is one of the polynomials $P_k$.
\end{lemma}
\begin{proof}
If $g$ has order $\ell$, the polynomial $X^\ell-1$ annihilates $g$
and is separable. Its nontrivial irreducible factors all have degree
six. A nonidentity module of dimension seven therefore consists of one
nontrivial irreducible constituent of dimension six and one trivial
constituent. There is no room for two nontrivial constituents. The
one-dimensional constituent is trivial because $\ell\nmid q-1$.
This gives $\chi_g=P_k$ for one coset.

Conversely $P_k$ divides $X^\ell-1$. If $\chi_g=P_k$, the
Cayley--Hamilton theorem gives $g^\ell=I$. The matrix is not the
identity, since $P_k\ne(X-1)^7$. Primality of $\ell$ proves its exact
order. In particular, no semisimplicity assumption is needed in this
direction: the annihilating polynomial already supplies it.
\end{proof}

\begin{lemma}\label{lem:distinct}
Every $P_k$ satisfies~\eqref{eq:polynomial}, and its signature
$\sigma_k$ is different from every other coset signature. Thus the
table in~\eqref{eq:table} has exactly $(\ell-1)/6$ entries.
\end{lemma}
\begin{proof}
In $\mathbb Z/\ell\mathbb Z$ the element $q$ has order six, so
$q^3=-1$ and $q\ne-1$. The factorization
$q^3+1=(q+1)(q^2-q+1)$ therefore gives $q^2=q-1$.
Its six powers are
\[
 1,\ q,\ q-1,\ -1,\ -q,\ 1-q.
\]
Taking $x=\zeta^k$, $y=\zeta^{k(q-1)}$ and
$z=\zeta^{-kq}$ in~\eqref{eq:weights} gives exactly the roots of
$P_k$. Consequently the same coefficient calculation applies to $P_k$.
If two signatures were equal, formula~\eqref{eq:polynomial} would make
their whole polynomials equal, contrary to the distinct Frobenius
orbits.
\end{proof}

\begin{theorem}\label{thm:table}
Let $g$ be an actual element of $G_2(q)$ in its acquired natural
seven-dimensional representation. Then
\[
 |g|=\ell\quad\Longleftrightarrow\quad
 (e_1(g),e_2(g))\in\{\sigma_k:kH\in(\mathbb Z/\ell\mathbb Z)^\times/H\}.
\]
\end{theorem}
\begin{proof}
Necessity follows from Lemma~\ref{lem:possible}. For sufficiency,
Lemma~\ref{lem:coefficients} and Lemma~\ref{lem:distinct} make equality
of signatures an equality of the entire characteristic polynomial.
Now apply the converse in Lemma~\ref{lem:possible}.
\end{proof}

For a general $\GL_7$ or orthogonal matrix, equality of two coefficients
does not prove equality of its characteristic polynomial. In that
setting a failed signature is a rejection test for prime order; a
matching signature by itself is inconclusive. Exact natural $G_2$
membership is the additional premise that makes the test sufficient.
Changing the basis of an already identified natural module changes
neither coefficient.

\section{Cost of obtaining the signature}

For a matrix $Y=(y_{ij})$, its trace uses only the seven diagonal
entries. Its second elementary coefficient is the sum of its principal
minors of order two:
\begin{equation}\label{eq:minors}
 e_2(Y)=\sum_{1\leq i<j\leq7}(y_{ii}y_{jj}-y_{ij}y_{ji}).
\end{equation}
This uses forty-two products of entries and no division. It is valid
in every characteristic, independently of any group membership.

In odd characteristic one may use the equivalent signature
$(t,u)=(\tr Y,\tr(Y^2))$. Newton's identity is
\[
 e_2(Y)=\frac{t^2-u}{2},
 \qquad
 u=\sum_{i=1}^7y_{ii}^2+2\sum_{1\leq i<j\leq7}y_{ij}y_{ji}.
\]
The last expression uses seven squares and twenty-one products of
opposite entries. Computing a whole matrix $Y^2$ is unnecessary. A
table stored directly as $(t,u)$ needs no conversion to $e_2$ at run
time. These are abstract field-operation counts, not measured Python
time, an acquired operation invoice or a new program. The same
coefficient calculation must not be used to accept a matrix without the
natural-module premise in Theorem~\ref{thm:table}.

\section{The order-$37$ specialization and a sixth-power filter}

At $q=11$, $\ell=37$ we have
\[
 H=\{1,11,10,36,26,27\}=\{\pm1,\pm10,\pm11\}.
\]
There are six distinct pairs. The present note does not compute their
field values or claim that a particular retained lookup table has been
authenticated. Nor can the natural $G_2$ premise remove any of the six
spectra. Once one order-$37$ element is present, its nonzero powers
realize all six Frobenius cosets. Existence follows, for example, from
the already imported unitary subgroup and Cauchy's theorem in the
retained prime-witness note. No new finite conjugacy-class count is
needed here.

We now reuse a different reviewed mathematical leg. The existing
involution theorem for a preserved alternating three-form with
nondegenerate contraction form forces every nonidentity involution to
have eigenspace dimensions $(3,4)$ for $(+1,-1)$. Its Cayley application
in characteristic eleven is retained in
\texttt{involutions\_preserving\_alternating\_forms.tex} and
\texttt{cayley\_contraction\_form.tex}; their scope and the exact
natural-model joins remain imported premises.

\begin{lemma}\label{lem:odd}
No nonidentity involution of natural $G_2(11)$ centralizes an element
of order $37$.
\end{lemma}
\begin{proof}
Let $Y$ have order $37$. The proof of Lemma~\ref{lem:possible} gives
the direct sum $V=L\oplus W$ of a trivial line and an irreducible
six-dimensional $\F_{11}[Y]$-module. A commuting linear operator
preserves both summands. The commutant on $W$ is the field
$\F_{11}[Y|_W]\cong\F_{11^6}$: $W$ is one-dimensional over that
field, and any commuting map is multiplication by its value on a
field basis vector.

If $z$ is a commuting involution, its restriction to $W$ is therefore
$I$ or $-I$, since these are the only roots of $X^2-1$ in this field.
Its restriction to $L$ is also $1$ or $-1$. Its negative eigenspace
has dimension $0,1,6$ or $7$. The first choice makes $z=I$; the others
contradict the previously reviewed multiplicity $(3,4)$ for a
nonidentity natural $G_2(11)$ involution.
\end{proof}

\begin{corollary}\label{cor:sixth}
For $w\in G_2(11)$ in its natural representation,
\[
 |w^6|=37\quad\Longleftrightarrow\quad |w|\in\{37,111\}.
\]
\end{corollary}
\begin{proof}
The forward assumption gives $w^{222}=I$ and $37\mid|w|$. Thus
$|w|$ is one of $37,74,111,222$. If it is even, the element
$w^{|w|/2}$ is a nonidentity involution centralizing $w^6$, contrary
to Lemma~\ref{lem:odd}. The two remaining orders give the reverse
implication by $|w^6|=|w|/\gcd(|w|,6)$.
\end{proof}

For order $111$, the Frobenius orbits on primitive roots again have
size six: $11^2\equiv10$ and $11^3\equiv-1\pmod{111}$. There are
$\varphi(111)/6=72/6=12$ such orbits. An order-$111$ matrix in
dimension seven has characteristic polynomial $(X-1)f$, with $f$ one
of these primitive degree-six factors. Indeed an order-$37$
constituent already has dimension six, so adjoining a separate
nontrivial order-three constituent would need two more dimensions.
The single remaining rational eigenvalue has order dividing both
$111$ and $10$, hence is one. It follows that a constituent of exact
order $111$ must occur.

Conversely each such characteristic polynomial implies exact order
$111$ by Cayley--Hamilton and the presence of a primitive eigenvalue.
Its six nonconstant eigenvalues have exponents
$k\{1,11,10,-1,-11,-10\}$, so the coefficient identity in
Lemma~\ref{lem:coefficients} applies. Its twelve signatures are
distinct from each other and from the six order-$37$ signatures.
Consequently an eighteen-pair table gives the following exact test:
\[
 |w^6|=37\quad\Longleftrightarrow\quad
 (\tr w,\tr(w^2))\text{ is one of those eighteen pairs},
 \qquad w\in G_2(11)\text{ naturally}.
\]
This establishes the required order reduction without asserting that a
full centralizer is cyclic of order $111$. It does not assert that any
new table or a saved word has been checked. An implementation may use
the statement only after exact table construction, Source review,
control execution and the acquired natural membership are supplied.

\section{Attribution and application boundary}

The natural weight description, finite-field factor degrees, elementary
coefficient identities, Newton's identity and Cayley--Hamilton are
standard mathematics. The coefficient reduction was independently
derived in the present discussion by \texttt{/root/forest\_resume}
and \texttt{/root/mathematical\_review}. No claim of historical novelty
is made. The sixth-power argument explicitly reuses the reviewed
involution and Cayley contraction lemmas; it is not a new classification
of finite centralizers.

The retained note
\texttt{computation/f004/theory/g2\_11\_generation\_by\_prime\_witnesses.tex}
supplies the order-$37$ module observation and the distinct generation
criterion involving a moved nonisotropic line. The spectral theorem
here proves orders only. It does not by itself prove generation, an
intersection property, a rank-four or rank-five statement, completion
of an Original domain, or any named conjecture. In particular the
geometric moved-line argument is not transferred to characteristic two.

There has been no numerical or symbolic program execution, factorization
run, Original operand activation, model or word evaluation, implementation
change, scientific case adoption, Lean formalization or compression
measurement in preparing this text. Applications require their own
exact membership, frame, Source, table, actual execution and independent
review joins.
\end{document}
